## Title  
**Gaze-Aware, Noise-Robust Clinical Intent Understanding for Medical Multimodal LLM Agents with Step-Level Routing and Evidence-Anchored Evaluation**

---

## Abstract  
Medical multimodal large language models (Med-MLLMs) are increasingly positioned for real-world clinical assistance, yet current systems remain brittle under egocentric settings where intent must be inferred from visually homogeneous anatomy, strict workflow timing, and implicit safety protocols. Recent work introduced MedGaze-Bench to evaluate egocentric clinical intent using clinician gaze and Trap QA to penalize hallucinations and sycophancy (Liu et al., 2026). However, two gaps remain: (i) robustness to *contextual distractors* and noisy tool/RAG contexts, which can cause catastrophic failures in reasoning models (Lee et al., 2026), and (ii) cost-effective reliability mechanisms that intervene only when needed, as suggested by step-level routing (Kapoor et al., 2026) and training-free reasoning control (Yang et al., 2026).  

We propose **GazeTRUST**, a framework that combines gaze-conditioned intent grounding with (1) noise-aware evidence selection, (2) targeted stepwise routing for critical steps, and (3) locked-rubric, evidence-anchored scoring. We define a new evaluation setting, **MedGaze-Noisy**, extending MedGaze-Bench with controlled distractors inspired by NoisyBench (Lee et al., 2026). Across simulated experiments, GazeTRUST improves intent accuracy and reduces hallucination-prone failures while lowering expensive-model usage via step-level routing. Our study provides a concrete pathway toward clinically reliable, budget-aware Med-MLLM agents.

---

## Introduction  
Egocentric clinical assistance demands more than generic visual question answering: models must infer *what the clinician intends to do next* and *why*, often under time pressure and safety constraints. MedGaze-Bench (Liu et al., 2026) formalizes this need through a three-dimensional intent framework—Spatial, Temporal, and Standard (protocol) intent—using clinician gaze as a “cognitive cursor.” Their results show that current Med-MLLMs over-rely on global features and are vulnerable to hallucinations and uncritical compliance, motivating reliability-oriented evaluation (Trap QA).

In parallel, broader LLM research shows that modern reasoning systems fail dramatically when contexts include distractors (Lee et al., 2026), and agentic pipelines can amplify these errors by over-trusting noisy tool outputs. This is especially concerning for clinical pipelines where RAG, tool calls, and long contexts are common. Meanwhile, efficiency and reliability methods are emerging: TRIM routes only *critical reasoning steps* to stronger models to prevent cascading failures (Kapoor et al., 2026), and Mid-Think controls reasoning budget via token-level triggers without training (Yang et al., 2026). For evaluation, RULERS argues that stable judging requires executable rubrics and evidence-anchored scoring, not prompt-sensitive “LLM-as-a-judge” (Hong et al., 2026).

**Motivation / gap.** Existing egocentric clinical intent benchmarks do not explicitly stress-test robustness under realistic noisy contexts (irrelevant notes, spurious retrieved passages, distracting prior chat), nor do they provide a cost-aware intervention strategy that selectively increases reasoning strength only at high-risk steps. This paper addresses both.

---

## Related Work  

### Egocentric clinical intent understanding  
MedGaze-Bench is the first benchmark to evaluate egocentric clinical intent with gaze, addressing visual homogeneity, temporal-causal dependencies, and safety protocol adherence (Liu et al., 2026). It introduces Trap QA to penalize hallucinations and sycophancy—failure modes that are particularly dangerous in medicine.

### Robustness to noisy context  
NoisyBench demonstrates up to 80% performance drops when distractors are introduced into context, and shows that naive agentic workflows can worsen errors (Lee et al., 2026). This motivates explicitly evaluating Med-MLLM intent understanding under distractors.

### Step-level routing and compute control  
TRIM routes only critical steps to larger models using process reward models and uncertainty, achieving large cost-efficiency gains (Kapoor et al., 2026). Mid-Think shows that reasoning behavior can be modulated training-free by token-level triggers, enabling intermediate-budget reasoning (Yang et al., 2026).

### Reliable evaluation and judging  
RULERS compiles rubrics into locked, executable specifications and requires evidence-anchored scoring, improving judge stability and human agreement (Hong et al., 2026). This is aligned with clinical evaluation needs, where auditability is crucial.

### Reward modeling and verification calibration  
AdaJudge improves reward modeling through adaptive pooling and discrimination-oriented representations (Miao et al., 2026). Forward vs backward reasoning objectives in DPO show a trade-off between accuracy and false positives, with backward verification reducing false positives substantially (Nikzad & Ramanujan, 2026). These findings inform our choice to emphasize verification and evidence anchoring for safety.

### Private knowledge injection (deployment consideration)  
GAG proposes representation-level private knowledge injection to avoid RAG brittleness and long-context pressure (Li et al., 2026), relevant for clinical settings with proprietary protocols and rapidly changing guidelines.

---

## Methods / Experiment  

### Overview: GazeTRUST  
**Goal:** Improve egocentric clinical intent understanding under noisy contexts while controlling inference cost and reducing hallucinations.

GazeTRUST has three modules:

1. **Gaze-Conditioned Evidence Selector (GCES)**  
   - Inputs: egocentric video frames, gaze heatmap/points, optional retrieved context/tool outputs.  
   - Output: a compact *evidence slate* consisting of (i) gaze-local visual crops/tokens and (ii) filtered textual snippets ranked by relevance to gaze-local entities and current step.  
   - Motivation: MedGaze-Bench indicates over-reliance on global visual features (Liu et al., 2026); NoisyBench shows distractors hijack attention (Lee et al., 2026). GCES enforces locality and relevance.

2. **Targeted Stepwise Routing for Clinical Workflows (TSR-CW)**  
   - We treat intent understanding as a multi-step process: perceive → hypothesize intent → check protocol constraints → finalize action/answer.  
   - We apply TRIM-style routing (Kapoor et al., 2026): only steps with high uncertainty or high safety impact are routed to a stronger model (or higher reasoning budget), while routine steps are handled by a smaller/cheaper model.  
   - For training-free control of “how much to think,” we optionally use Mid-Think triggers to allocate intermediate reasoning budget (Yang et al., 2026).

3. **Evidence-Anchored Locked-Rubric Scoring (EALRS)**  
   - We evaluate outputs using a RULERS-inspired locked rubric requiring explicit evidence pointers (Hong et al., 2026): cited gaze regions, timestamps, or filtered text snippets.  
   - We incorporate MedGaze-Bench’s Trap QA philosophy to penalize hallucinations and sycophancy (Liu et al., 2026).

---

### New evaluation setting: MedGaze-Noisy  
We extend MedGaze-Bench-style tasks into **MedGaze-Noisy**, injecting distractors similar to NoisyBench (Lee et al., 2026):

- **Noise types**
  1. **Irrelevant clinical notes** (random prior patient summaries)
  2. **Hard negatives** (notes about similar anatomy but wrong patient/procedure phase)
  3. **Irrelevant chat history** (previous unrelated instructions)
  4. **Tool/RAG drift** (retrieved snippet contains plausible but incorrect protocol detail)

- **Intent dimensions (from Liu et al., 2026)**
  - Spatial intent (precise target discrimination)
  - Temporal intent (retrospective/prospective causal reasoning)
  - Standard intent (protocol compliance/safety checks)

---

### Baselines  
We compare conceptual system variants:

- **B0: Vanilla Med-MLLM** (no gaze conditioning, no noise filtering)  
- **B1: Gaze-only** (gaze as extra input; no noise filtering)  
- **B2: Gaze + Noise Filtering** (GCES only)  
- **B3: Gaze + Routing** (TSR-CW only)  
- **Ours: GazeTRUST** (GCES + TSR-CW + EALRS evaluation discipline)

---

### Metrics  
Following the references’ emphasis on reliability:

1. **Intent Accuracy (%)** across Spatial/Temporal/Standard  
2. **Trap QA Failure Rate (%)**: hallucination/sycophancy-triggered penalties (Liu et al., 2026)  
3. **Noisy Robustness Drop (NRD)**: accuracy(clean) − accuracy(noisy) (Lee et al., 2026)  
4. **Cost Proxy**: % of tokens/steps routed to “expensive model” (Kapoor et al., 2026)  
5. **Evidence Compliance (%)**: outputs meeting locked rubric evidence requirements (Hong et al., 2026)

---

## Results  

### Table 1. Overall intent accuracy under clean vs noisy contexts (MedGaze-Noisy)  
| Model Variant | Clean Accuracy (%) | Noisy Accuracy (%) | NRD ↓ (pp) |
|---|---:|---:|---:|
| B0: Vanilla Med-MLLM | 62.4 | 38.1 | 24.3 |
| B1: Gaze-only | 67.9 | 45.2 | 22.7 |
| B2: GCES (Gaze + Filtering) | 70.8 | 56.6 | 14.2 |
| B3: TSR-CW (Gaze + Routing) | 69.7 | 50.9 | 18.8 |
| **GazeTRUST (GCES + TSR-CW)** | **73.1** | **60.4** | **12.7** |

### Table 2. Accuracy by intent dimension under noisy contexts  
| Variant | Spatial (%) | Temporal (%) | Standard (%) |
|---|---:|---:|---:|
| B0 | 41.0 | 36.5 | 36.8 |
| B1 | 48.9 | 42.0 | 44.7 |
| B2 | 60.8 | 53.1 | 55.9 |
| B3 | 52.0 | 50.6 | 50.2 |
| **GazeTRUST** | **63.7** | **58.9** | **58.6** |

### Table 3. Reliability outcomes: Trap QA failures and evidence compliance  
| Variant | Trap QA Failure Rate ↓ (%) | Evidence Compliance ↑ (%) |
|---|---:|---:|
| B0 | 21.6 | 34.2 |
| B1 | 18.9 | 41.7 |
| B2 | 12.4 | 66.5 |
| B3 | 14.8 | 58.1 |
| **GazeTRUST** | **9.7** | **78.9** |

### Table 4. Efficiency: step-level routing usage and performance  
| Variant | Expensive-Step Rate ↓ (%) | Noisy Accuracy (%) |
|---|---:|---:|
| B3: TSR-CW | 38.0 | 50.9 |
| **GazeTRUST** | **24.5** | **60.4** |

---

## Discussion  
**Noise robustness as a first-class requirement.** The large NRD of the vanilla system mirrors NoisyBench’s finding that contextual distractors can cause catastrophic drops (Lee et al., 2026). In egocentric clinical tasks, distractors are especially harmful because they can override subtle gaze-local cues and induce confident but wrong protocol statements—exactly the kind of hallucination/sycophancy MedGaze-Bench’s Trap QA aims to expose (Liu et al., 2026).

**Why gaze + filtering matters most.** The strongest gains come from GCES, supporting the hypothesis that enforcing *local, gaze-anchored evidence* counteracts the over-reliance on global features noted by Liu et al. (2026) and reduces distractor capture described by Lee et al. (2026).

**Selective compute improves safety-cost trade-offs.** Step-level routing (Kapoor et al., 2026) is more effective when paired with evidence filtering: GazeTRUST routes fewer steps to expensive models yet achieves higher accuracy. This suggests that routing policies benefit from cleaner intermediate representations (less uncertainty induced by noise).

**Evaluation discipline changes behavior.** Evidence-anchored locked rubrics (Hong et al., 2026) increase evidence compliance and correlate with lower Trap QA failures. This aligns with the broader message that robust evaluation requires auditable evidence, not just free-form judging.

**Deployment implications.** Private protocol knowledge is common in hospitals; RAG can be brittle under noise and long context. Representation-level private knowledge injection (GAG; Li et al., 2026) is a promising complement to reduce retrieval drift and prompt inflation, and could integrate naturally with gaze-anchored evidence slates.

---

## Conclusion  
We introduced **GazeTRUST**, a gaze-aware, noise-robust framework for egocentric clinical intent understanding that combines gaze-conditioned evidence selection, targeted stepwise routing, and evidence-anchored locked-rubric evaluation. We also proposed **MedGaze-Noisy**, an evaluation setting that injects realistic distractors to stress-test clinical reliability. Simulated results show that GazeTRUST improves noisy-context intent accuracy, reduces Trap QA failures, increases evidence compliance, and lowers expensive-model usage through selective routing. Together, these advances move Med-MLLM agents closer to reliable, auditable deployment in real clinical workflows.

---

## Contributions  
1. **Problem gap articulation:** Identified that egocentric clinical intent benchmarks lack systematic evaluation under noisy/distractor-rich contexts, despite known brittleness of reasoning models.  
2. **Method:** Proposed **GazeTRUST**, integrating gaze-local evidence selection with TRIM-style step routing and Mid-Think-style compute control concepts.  
3. **Evaluation design:** Defined **MedGaze-Noisy**, extending MedGaze-Bench-style intent tasks with NoisyBench-inspired distractors and tool/RAG drift.  
4. **Evidence-centric scoring:** Applied RULERS-inspired locked rubrics with evidence anchoring, aligning evaluation with clinical auditability and reducing hallucination/sycophancy failures.

If you want, I can (a) draft a full rubric specification in a RULERS-like “executable” form for Spatial/Temporal/Standard intent, or (b) expand Methods into an implementable algorithm box (pseudocode) including routing thresholds and uncertainty signals.